% from ega1-3-fr.tex
\oldpage[I]{117}
\phantomsection\label{I.3.5.10.base-change-points-diagram-fr}
\[
  \xymatrix{
    X'(K)\ar[r]\ar[d] & Y'(K)\ar[d]\\
    X(K)\ar[r] & Y(K)
  }
\]

\begin{remark}[3.5.11]
\label{I.3.5.11-fr}
Nous dirons qu'un morphisme $f=(\psi,\theta)$ de préschémas est
\emph{injectif} si l'application $\psi$ est injective. Pour qu'un morphisme
$f=(\psi,\theta):X\to Y$ soit radiciel, il faut et il suffit que pour tout
morphisme $Y'\to Y$, le morphisme $f_{(Y')}:X_{(Y')}\to Y'$ soit injectif (ce
qui justifie la terminologie de morphisme \emph{universellement injectif}). En
effet, la condition est nécessaire en vertu de
(\hyperref[I.3.5.7-fr]{3.5.7, (ii)}) et de
(\hyperref[I.3.5.8-fr]{3.5.8}). Inversement, la condition implique d'abord que
$\psi$ est injectif~; si pour un $x\in X$, le monomorphisme
$\theta^x:k(\psi(x))\to k(x)$ n'était pas radiciel, il y aurait une extension
$K$ de $k(\psi(x))$ et deux morphismes distincts
$\operatorname{Spec}(K)\to X$ correspondant au même morphisme
$\operatorname{Spec}(K)\to Y$ (\hyperref[I.3.5.8-fr]{3.5.8}). Mais alors, en
posant $Y'=\operatorname{Spec}(K)$, il y aurait deux $Y'$-sections distinctes
de $X_{(Y')}$ (\hyperref[I.3.3.14-fr]{3.3.14}), ce qui contredit l'hypothèse
que $f_{(Y')}$ est injectif.
\end{remark}

\subsection{Fibres.}
\label{subsection:I.3.6-fr}

\begin{proposition}[3.6.1]
\label{I.3.6.1-fr}
Soient $f:X\to Y$ un morphisme, $y$ un point de $Y$, $\mathfrak{a}_y$ un
idéal de définition de $\mathcal{O}_y$ pour la topologie
$\mathfrak{m}_y$-préadique. La projection
$p:X\times_Y\operatorname{Spec}(\mathcal{O}_y/\mathfrak{a}_y)\to X$ est un
homéomorphisme de l'espace sous-jacent au préschéma
$X\times_Y\operatorname{Spec}(\mathcal{O}_y/\mathfrak{a}_y)$ sur la fibre
$f^{-1}(y)$ munie de la topologie induite par celle de l'espace sous-jacent à
$X$.
\end{proposition}

Comme $\operatorname{Spec}(\mathcal{O}_y/\mathfrak{a}_y)\to Y$ est radiciel
(\hyperref[I.3.5.4-fr]{3.5.4} et \hyperref[I.2.4.7-fr]{2.4.7}) et que
$\operatorname{Spec}(\mathcal{O}_y/\mathfrak{a}_y)$ est réduit à un seul
point, l'idéal $\mathfrak{m}_y/\mathfrak{a}_y$ étant nilpotent par hypothèse
(\hyperref[I.1.1.12-fr]{1.1.12}), on sait déjà
(\hyperref[I.3.5.10-fr]{3.5.10} et \hyperref[I.3.3.4-fr]{3.3.4}) que $p$
identifie en tant \emph{qu'ensemble} l'espace sous-jacent à
$X\times_Y\operatorname{Spec}(\mathcal{O}_y/\mathfrak{a}_y)$ à $f^{-1}(y)$~;
tout revient à démontrer que $p$ est un homéomorphisme. En vertu de
(\hyperref[I.3.2.7-fr]{3.2.7}), la question est locale sur $X$ et $Y$, et on
peut donc supposer que $X=\operatorname{Spec}(B)$,
$Y=\operatorname{Spec}(A)$, $B$ étant une $A$-algèbre. Le morphisme $p$
correspond alors à l'homomorphisme
$1\otimes\varphi:B\to B\otimes_A A'$, où $A'=A_y/\mathfrak{a}_y$ et
$\varphi$ est l'application canonique de $A$ dans $A'$. Or, tout élément de
$B\otimes_A A'$ s'écrit
\phantomsection\label{I.3.6.1.localization-fraction-fr}
\[
  \sum_i b_i\otimes\varphi(a_i)/\varphi(s)
  =\left(\sum_i(a_i b_i\otimes 1)\right)(1\otimes\varphi(s))^{-1},
\]
où $s\notin\mathfrak{j}_y$, et la prop.
(\hyperref[I.1.2.4-fr]{1.2.4}) est applicable.

\begin{env}[3.6.2]
\label{I.3.6.2-fr}
Dans toute la suite de ce Traité, lorsque nous considérerons une fibre
$f^{-1}(y)$ d'un morphisme comme munie d'une structure de $k(y)$-préschéma,
\emph{il s'agira toujours du préschéma obtenu en transportant la structure de
$X\times_Y\operatorname{Spec}(k(y))$ par la projection dans $X$}. Nous
écrirons aussi ce dernier produit $X\otimes_Y k(y)$, ou
$X\otimes_{\mathcal{O}_Y}k(y)$~; plus généralement, si $B$ est une
$\mathcal{O}_y$-algèbre, nous noterons $X\otimes_Y B$ ou
$X\otimes_{\mathcal{O}_Y}B$ le produit
$X\times_Y\operatorname{Spec}(B)$.
\end{env}

Avec la convention précédente, il résulte de
(\hyperref[I.3.5.10-fr]{3.5.10}) que les points de $X$ à valeurs dans une
extension $K$ de $k(y)$ sont identifiés aux \emph{points de $f^{-1}(y)$ à
valeurs dans $K$}.

\begin{env}[3.6.3]
\label{I.3.6.3-fr}
Soient $f:X\to Y$, $g:Y\to Z$ deux morphismes, $h=g\circ f$ leur composé~;
pour tout $z\in Z$, la fibre $h^{-1}(z)$ est un préschéma isomorphe à
\phantomsection\label{I.3.6.3.fibre-composition-fr}
\[
  X\times_Z\operatorname{Spec}(k(z))
  =(X\times_Y Y)\times_Z\operatorname{Spec}(k(z))
  =X\times_Y g^{-1}(z)
\]
